% Requires \usepackage{booktabs}
\begin{table}[htbp]
\centering
\small
\caption{Country-crop systems with high research need and no eligible study, ranked by share of world production.}
\label{tab:tab_03_under_researched_systems}
\begin{tabular}{llllrrrrrrr}
\toprule
ISO3 & Country & Crop & World Bank region & Income group & Production (Mt) & Harvested area (Mha) & Share of world production (\%) & Need percentile & Fractional studies & Full-count studies \\
\midrule
NGA & Nigeria & yam & Sub-Saharan Africa  & Lower middle income & 57.12 & 6.86 & 0.782 & 92.2 & 0 & 0 \\
COD & Democratic Republic of the Congo & cassava & Sub-Saharan Africa  & Low income & 41.61 & 5.11 & 0.570 & 94.8 & 0 & 0 \\
GTM & Guatemala & sugarcane & Latin America \& Caribbean  & Upper middle income & 27.31 & 0.26 & 0.374 & 75.6 & 0 & 0 \\
MMR & Myanmar & sugarcane & East Asia \& Pacific & Lower middle income & 11.40 & 0.18 & 0.156 & 78.8 & 0 & 0 \\
NGA & Nigeria & oilpalm & Sub-Saharan Africa  & Lower middle income & 10.20 & 3.97 & 0.140 & 92.2 & 0 & 0 \\
UGA & Uganda & banana plantain & Sub-Saharan Africa  & Low income & 10.01 & 2.24 & 0.137 & 97.4 & 0 & 0 \\
BOL & Bolivia (Plurinational State of) & sugarcane & Latin America \& Caribbean  & Lower middle income & 9.80 & 0.18 & 0.134 & 75.1 & 0 & 0 \\
AGO & Angola & cassava & Sub-Saharan Africa  & Lower middle income & 9.55 & 0.98 & 0.131 & 93.3 & 0 & 0 \\
ECU & Ecuador & sugarcane & Latin America \& Caribbean  & Upper middle income & 9.38 & 0.12 & 0.128 & 80.8 & 0 & 0 \\
TZA & United Republic of Tanzania & cassava & Sub-Saharan Africa  & Lower middle income & 9.01 & 1.27 & 0.123 & 95.9 & 0 & 0 \\
IRN & Iran (Islamic Republic of) & sugarcane & Middle East, North Africa, Afghanistan \& Pakistan & Upper middle income & 8.41 & 0.09 & 0.115 & 76.2 & 0 & 0 \\
CIV & Côte d'Ivoire & yam & Sub-Saharan Africa  & Lower middle income & 7.57 & 1.36 & 0.104 & 81.9 & 0 & 0 \\
MWI & Malawi & sweet potato & Sub-Saharan Africa  & Low income & 6.78 & 0.30 & 0.093 & 91.2 & 0 & 0 \\
MOZ & Mozambique & cassava & Sub-Saharan Africa  & Low income & 6.22 & 0.79 & 0.085 & 97.9 & 0 & 0 \\
CIV & Côte d'Ivoire & cassava & Sub-Saharan Africa  & Lower middle income & 6.21 & 1.09 & 0.085 & 81.9 & 0 & 0 \\
NGA & Nigeria & banana plantain & Sub-Saharan Africa  & Lower middle income & 5.89 & 0.49 & 0.081 & 92.2 & 0 & 0 \\
CMR & Cameroon & cassava & Sub-Saharan Africa  & Lower middle income & 5.86 & 0.42 & 0.080 & 78.2 & 0 & 0 \\
MWI & Malawi & cassava & Sub-Saharan Africa  & Low income & 5.85 & 0.24 & 0.080 & 91.2 & 0 & 0 \\
ZWE & Zimbabwe & sugarcane & Sub-Saharan Africa  & Lower middle income & 5.76 & 0.07 & 0.079 & 89.1 & 0 & 0 \\
COD & Democratic Republic of the Congo & banana plantain & Sub-Saharan Africa  & Low income & 5.67 & 1.30 & 0.078 & 94.8 & 0 & 0 \\
\bottomrule
\end{tabular}
\begin{minipage}{\linewidth}\footnotesize
Selection rule, fixed before the table was produced: country-crop cells with a fractional eligible study count of exactly zero, a need percentile at or above the 75th, and a positive FAOSTAT production denominator, ranked by share of world production. 643 of the 2,616 panel cells meet the rule, spanning 49 countries; the 20 largest by production share are shown. A zero here is measured, not missing: the cell was in the panel and no eligible study covered it. The need percentile is the rank-weighted need index rescaled from 0-1 to 0-100 and is a country-level quantity, identical for every crop within a country. The full ranking is in the supplement.
\end{minipage}
\end{table}
